\subsubsection{应力-能量张量} % 原文: 4.2.1 Stress-Energy Tensor

一般而言，与引力的耦合定义了能量-动量张量（energy-momentum tensor）：

\begin{equation}
  T_{\mu\nu}:=\frac{1}{2\sqrt{g}}\frac{\delta}{\delta g^{\mu\nu}}S
  \tag{4.74}
  \label{eq:4.74}
\end{equation}

% [Tutor Note: (4.74) 是「对度规变分」的定义式，原文既未对 $\mu\nu$ 作对称化，也没区分 $\delta/\delta g^{\mu\nu}$ 与 $\delta/\delta g_{\mu\nu}$。核对 (4.76) 时极易出错的地方正在这里：那里给出的是对 $\delta h^{--}$ 的变分，而 $\sqrt{g}$ 本身也依赖度规，因此 $\delta S/\delta h^{--}$ 与 $T_{--}$ 之间要先按 $\delta g^{\mu\nu}$ 的换算关系过渡，不能直接等同。]

让我们继续在光锥坐标（light-cone coordinates）下工作。经典共形不变性表明 $T_{+-}=0$。在变分 $\delta h^{--}$ 之下我们有：

\begin{equation}
  \begin{aligned}
    \partial_- &\to \partial_-\\
    \partial_+ &\to \partial_+ + \delta h^{--}\partial_-
  \end{aligned}
  \tag{4.75}
  \label{eq:4.75}
\end{equation}

由于 WZ 项不依赖于度规，我们容易算出

\begin{equation}
  \begin{aligned}
    \delta S &= -6\pi c_1 k\operatorname{Tr}(g^{-1}\partial_- g)^2\,\delta h^{--}\\
    T_{--} &= \frac{\pi}{k}\operatorname{Tr}(J_-)^2
  \end{aligned}
  \tag{4.76}
  \label{eq:4.76}
\end{equation}

而 $T_{++}$ 完全类似。

哈密顿量为

\begin{equation}
  H\sim\oint dx\,(T_{++}+T_{--})
  \tag{4.77}
  \label{eq:4.77}
\end{equation}

把它分解为 Fourier 模式是方便的

\begin{equation}
  T_{--}(x)=\mathit{const.}\sum_{n\in\mathbb{Z}}L_n e^{2\pi i n(t-x)}
  \tag{4.78}
  \label{eq:4.78}
\end{equation}

\begin{equation}
  T_{++}(x)=\mathit{const.}\sum_{n\in\mathbb{Z}}\tilde L_n e^{2\pi i n(t+x)}
  \tag{4.79}
  \label{eq:4.79}
\end{equation}

其中常数取得使得

\begin{equation}
  L_n=\frac{1}{2k}\sum_{m\in\mathbb{Z}}J^a_{n-m}J^a_m
  \tag{4.80}
  \label{eq:4.80}
\end{equation}

在量子理论中，我们将定义一个「最高权」真空，使得 $J^a_n|vac\rangle=0$（$n>0$）。因此，我们需要对表达式 \eqref{eq:4.80} 作正规序化（normal-ordering），即把 $n>0$ 的 $J^a_n$ 移到右边。

由于当 $x$ 接近 $y$ 时 $J(x)$ 与 $J(y)$ 相乘会出现短距离奇异性，我们应当允许用两个流之积表示的这个表达式带有乘性重整化（multiplicative renormalization）。

现在，$L_0$ 生成标度变换，而 $J^a_-$ 是守恒流。因此，在量子理论中 $J^a_-$ 的标度维数（scaling dimension）必须恰好等于一。这意味着必须有

\begin{equation}
  [L_0,J^a_n]=-nJ^a_n
  \tag{4.81}
  \label{eq:4.81}
\end{equation}

另一方面，我们可以算出：

\begin{equation}
  \Big[\sum_{m\in\mathbb{Z}}:J^b_{-m}J^b_m:,J^a_n\Big]=-2n(k+h^\vee)J^a_n
  \tag{4.82}
  \label{eq:4.82}
\end{equation}

其中我们用到了这样一个事实：在正交归一（ON）基下

\begin{equation}
  \sum_{bc}f^{abc}f^{dbc}=2h^\vee\delta^{ad}
  \tag{4.83}
  \label{eq:4.83}
\end{equation}

其中 $h$ 是对偶 Coxeter 数（dual Coxeter number）。

% [Tutor Note: 原文此处写作「$h$ 是对偶 Coxeter 数」，紧接其前的 (4.83) 与紧接其后的 (4.84) 里却都带着上标 $\vee$。按通常记号应统一为 $h^\vee$（即对偶 Coxeter 数，$\mathfrak{su}(N)$ 时为 $N$）。照原文转写为 $h$，引用时请自行补上 $\vee$。]

因此，我们得知量子效应导致了如下重整化

\begin{equation}
  \frac{1}{k}\to\frac{1}{k+h^\vee}
  \tag{4.84}
  \label{eq:4.84}
\end{equation}

回忆一下，半经典展开参数是 $1/k$。因此，只要 $k+h\neq0$，量子应力-能量张量就是\footnote{当 $k+h^\vee=0$ 时，存在与流代数对易的新算符，对称性甚至更大。这一情形与几何 Langlands 纲领（Geometric Langlands Program）有关。}

\begin{equation}
  L_n=\frac{1}{2(k+h^\vee)}\sum_{m\in\mathbb{Z}}:J^a_mJ^a_{n-m}:
  \tag{4.85}
  \label{eq:4.85}
\end{equation}

现在人们可以直接（尽管有点繁琐地）算出

\begin{equation}
  [L_n,L_m]=(n-m)L_{n+m}+\frac{c}{12}(n^3-n)\delta_{n+m,0}
  \tag{4.86}
  \label{eq:4.86}
\end{equation}

其中

\begin{equation}
  c=\frac{k}{k+h^\vee}\dim\mathfrak{g}
  \tag{4.87}
  \label{eq:4.87}
\end{equation}

注意当 $k\to\infty$ 时它趋于 $\dim\mathfrak{g}$。

\textbf{评注：}

\begin{itemize}
  \item 既然我们可以在 Hilbert 空间上表示 Virasoro 代数，我们就得到了一个共形场论。因此，在量子理论中可以维持 $T_{+-}=0$。耦合常数的取值 $f^2=\pm12\pi c_1 k$ 定义了理论 beta 函数的一个精确零点。若 $f^2$ 不取此值，模型是渐近自由的（asymptotically free）。
  \item 由于中心荷非零，理论的左行与右行部分具有\emph{引力反常}（gravitational anomaly）。正如理论经典的 $L\mathfrak{g}$ 对称性到头来是反常的、在量子理论中被 $\widehat{L\mathfrak{g}}$ 取代一样，理论经典的 $\mathrm{diff}(S^1)$ 对称性也被 Virasoro 代数取代。
  \item 尽管左行与右行部分各自带有反常，圆的微分同胚的生成元是 $L_n-\tilde L_{-n}$，而它们实际上没有反常。因此，我们可以自洽地把模型与引力耦合。
  \item 上面的对易子可以直接算出，但有一个更优雅的形式体系，它基于全纯场的算符乘积展开（operator product expansion），能给出同样的答案。

  % [Tutor Note: 作者说对易子「可以直接（尽管有点繁琐地）」算出，被略去的是对 (4.85) 中正规序乘积作逐次 Wick 收缩、并用 (4.83) 把 $f^{abc}f^{dbc}$ 收缩成 $2h^\vee\delta^{ad}$ 的那一步；等价地可以用 $J^a(z)J^b(w)\sim\frac{k\delta^{ab}}{(z-w)^2}+\frac{if^{abc}J^c(w)}{z-w}$ 的 OPE 一次得到 (4.86) 的 $c$。作者在最后一条评注里已明说 OPE 是更优雅的路线。]

  \item 中心荷可以解释为对非线性 sigma 模型目标空间维数的一种度量。若 $r$ 是该单 Lie 代数的秩，则
  \begin{equation}
    r\le c\le\dim\mathfrak{g}\qquad 1\le k<\infty
    \tag{4.88}
    \label{eq:4.88}
  \end{equation}

  在半经典极限 $k\to\infty$ 下，目标空间的维数是 $\dim\mathfrak{g}$：模型的零模部分就已经成为粒子在群流形上的量子力学。然而，对有限的 $k$ 必须使用某种「量子几何」的概念。这在 $k=1$ 时尤为戏剧化。此时对单边（simply laced）的 $ADE$ 群有 $c=r$。\footnote{对所有单群，$d/(g+1)=r$，其中 $g$ 是 Coxeter 数。} 此时事实上存在一个等价的 sigma 模型，其目标空间就是 Cartan 环面！这是「非阿贝尔玻色化（nonabelian bosonization）」与 Frenkel-Kac-Segal 构造的一部分。

  \item $L_0$ 的 Sugawara 构造清楚地表明：它就是 $L\mathfrak{g}$ 的二次 Casimir 算符。换一种说法，它是无穷维圈群上的 Laplacian。
\end{itemize}

\medskip\noindent\textbf{练习}\quad 证明

\begin{equation}
  [L_n,J^a_m]=-mJ^a_{n+m}
  \tag{4.89}
  \label{eq:4.89}
\end{equation}

\subsection{圈群的中心扩张与 WZ 项} % 原文: 4.3 The Central Extension Of The Loop Group And The WZ Term

\begin{keytakeaways}
  \item 几何构造：圈群 $LG$ 的中心扩张不必直接造在 $LG$ 上，而可以先在「圆盘上的映射群」$DG=\mathrm{Map}(D,G)$ 上写出显式的上闭链 (4.90)--(4.91)——那里有圆盘 $D$ 可供积分。
  \item 关键恒等式：Wess-Zumino 泛函满足 $WZ(g_1g_2)=WZ(g_1)+WZ(g_2)+6\pi kc_1\int_D\operatorname{Tr}(dg_2g_2^{-1})(g_1^{-1}dg_1)$，即该上闭链限制到「边界上为恒等」的子群 $D_1G$ 时变成上边界，于是可以除以 $\psi(D_1G)$ 得到中心扩张的圈群 $\widehat{LG}$。
  \item 与三维 Chern-Simons 的联系：$D_1G$ 正是 $D\times\mathbb{R}$ 上 Chern-Simons 理论的规范变换群，圆盘上的平坦规范场由 $DG$ 参数化，因此这个表述对 CS 理论的量子化极为方便。
  \item 附加的 $L_0$：计入 $L_0$ 后它生成圈的刚性旋转，$\exp[i\theta_0L_0]g(\theta)\exp[-i\theta_0L_0]=g(\theta+\theta_0)$，这为再与 Virasoro 群作半直积铺路。
\end{keytakeaways}

我们将给出圈群 $LG$ 的上闭链（cocycle）的一个非常漂亮的几何构造，这个上闭链定义了中心扩张的圈群。

窍门是考虑从\emph{圆盘}到群 $G$ 的映射构成的群 $DG$，也就是说，我们引入 $DG=\mathrm{Map}(D,G)$，其中 $D$ 是圆盘。注意，满足 $g|_{\partial D}=1$ 的映射构成的子群 $D_1G$ 是一个正规子群，并且 $DG/D_1G\cong LG$，显式的同构由限制映射给出。

现在，与 $LG$ 相反，写出群 $DG$ 的一个中心扩张 $\widehat{DG}$ 是容易的：

\begin{equation}
  (g_1,\lambda_1)\cdot(g_2,\lambda_2)=(g_1g_2,\lambda_1\lambda_2f(g_1,g_2))\qquad g_i\in DG
  \tag{4.90}
  \label{eq:4.90}
\end{equation}

其中

\begin{equation}
  f(g_1,g_2)=\exp\left[2\pi i(6\pi c_1k)\int_D\operatorname{Tr}(dg_2g_2^{-1})(g_1^{-1}dg_1)\right]
  \tag{4.91}
  \label{eq:4.91}
\end{equation}

注意我们把 Ad 不变的内积 $(\cdot,\cdot)_\mathfrak{g}$ 写成了某个表示中确定的迹 $\operatorname{Tr}$。对 $SU(N)$，取 $N$ 维表示中的迹时 $c_1=1/(24\pi^2)$。

\medskip\noindent\textbf{练习}\quad

a.) 验证 \eqref{eq:4.91} 确实是一个群上闭链。

b.) 计算 Lie 代数 $D\mathfrak{g}$ 上相应的中心扩张，并证明当其中一个元素在边界上为零时它是平凡的。具体地，证明它为

\begin{equation}
  24\pi^2ic_1k\oint_{S^1}\operatorname{Tr}\epsilon_1d\epsilon_2=ik\oint\operatorname{Tr}\epsilon_1d\epsilon_2
  \tag{4.92}
  \label{eq:4.92}
\end{equation}

其中 $c_1=1/(24\pi^2)$。

现在，美妙的观察是：当 $g_1$ 与 $g_2$ 在边界 $\partial D$ 上等于 1 时，我们可以把它们看作定义了从 $S^2$ 到 $G$ 的映射，从而可以定义 WZ 项。但是，由于我们上面证明的恒等式：

\begin{equation}
  WZ(g_1g_2)=WZ(g_1)+WZ(g_2)+6\pi kc_1\int\operatorname{Tr}(dg_2g_2^{-1})(g_1^{-1}dg_1)
  \tag{4.93}
  \label{eq:4.93}
\end{equation}

该上闭链限制到子群 $D_1G$ 上时就变成了上边界。因此，扩张

\begin{equation}
  1\to U(1)\to\widehat{DG}\to DG\to1
  \tag{4.94}
  \label{eq:4.94}
\end{equation}

在正规子群 $D_1G$ 上分裂，也就是说：

\begin{equation}
  \psi:g\to(g,e^{iWZ(g)})\qquad g\in D_1G
  \tag{4.95}
  \label{eq:4.95}
\end{equation}

是从 $D_1G$ 到 $\widehat{DG}$ 的群同态；于是我们可以取商

\begin{equation}
  1\to U(1)\to\widehat{DG}/\psi(D_1G)\to DG/D_1G=LG\to1
  \tag{4.96}
  \label{eq:4.96}
\end{equation}

来构造圈群 $\widehat{LG}:=\widehat{DG}/\psi(D_1G)$。

最后，如果计入 $L_0$，则注意

\begin{equation}
  \exp[i\theta_0L_0]g(\theta)\exp[-i\theta_0L_0]=g(\theta+\theta_0)
  \tag{4.97}
  \label{eq:4.97}
\end{equation}

所以 $L_0$ 生成圈的刚性旋转。

\textbf{评注}

\begin{enumerate}
  \item 上述中心扩张的构造归功于 J. Mickelsson。关于对任意流形 $X$ 的 $\mathrm{Map}(X,G)$ 的推广，见 [37] 及其中的参考文献。
  \item 在 Lie 群层面，人们可以与 Virasoro 群——圆的中心扩张微分同胚群——构造半直积。
  \item 上述中心扩张圈群的表述，对于在 $D\times\mathbb{R}$ 上量子化三维 Chern-Simons 理论非常方便。规范变换群是 $D_1G$。圆盘上的平坦规范场由 $DG$ 参数化。
\end{enumerate}

% [Tutor Note: 从 (4.93) 到「上闭链变上边界」少了一步说明：约束到 $D_1G$ 后 $D$ 上的积分可写成 $S^2$ 上的 WZ 泛函差，故 $f(g_1,g_2)$ 恰好等于 $e^{iWZ}$ 的余边界；这一步之所以成立，依赖 WZ 项的 $\bmod\,2\pi$ 整性（即 $c_1$ 的选取），也正是 (4.92) 那道练习要证的内容。]

\subsection{与反常的关系} % 原文: 4.4 Relation To Anomalies

\begin{keytakeaways}
  \item 主线：WZ 项不是凭空添加的拓扑项，它是量子有效作用量的反常变分之积分；二维 WZW 模型因此可以看作「把反常写成经典作用量」的产物。
  \item 具体机制：把 $\frac{\delta}{\delta A_z}\Psi$ 认作流的一点函数 $\langle J^{0,1}(x)\rangle$，WZW 的有效作用量就满足反常方程 (4.99)。
  \item 费米子计算：一个复手征费米子的单圈图给出 $\Gamma(A_-)\sim\int d^2q\,A_-(q)A_-(-q)\frac{q_+}{q_-}$，其规范变分 $\delta\Gamma\sim\int d^2x\,F(x)\epsilon(x)$ 就是二维 $U(1)$ 反常；非阿贝尔情形多出与 $A_-$ 无关的 $F(g)$，它正是 WZ 项。
  \item 结构结论：二维中可写 $A_-=h^{-1}\partial_-h$，于是 $\Gamma(A_-)\sim iS^+_{WZW}(h)$——量子反常的效应被一个经典作用量完整概括。
\end{keytakeaways}

\subsubsection{与外部规范场耦合} % 原文: 4.4.1 Coupling To External Gauge Fields

当我们把整体对称性 $LG_{hol}\times LG_{anti-hol}$ 规范化时，我们得到规范化的 WZW 作用量 (3.68)

{\small\itshape ［原文标记：或者是涉及 $S_+$ 的那一个］}
% [原文待办标记（页边），原文 p.118]

为简单起见取 $A^{0,1}=0$。于是，把 WZW 理论中的有效作用量等同于 Chern-Simons 波函数

\begin{equation}
  Z=\Psi[A_z]
  \tag{4.98}
  \label{eq:4.98}
\end{equation}

并把 $\frac{\delta}{\delta A_z}\Psi$ 等同于流的一点函数 $\langle J^{0,1}(x)\rangle$，我们就可以把 (3.47) 读作一个反常方程：

\begin{equation}
  D_{A,z}\langle J_{\bar z}\rangle=-\frac{k}{\pi}\bar\partial_{\bar z}A_z
  \tag{4.99}
  \label{eq:4.99}
\end{equation}

% [Tutor Note: 原文此句语序残缺（「... we can (3.47) as an anomaly equation」缺动词），按上下文只能理解为「把 (3.47) 读作反常方程」。另注意 (4.99) 左边是协变导数 $D_{A,z}$、右边是 $\bar\partial_{\bar z}A_z$，而 (4.98) 把 $Z$ 等同于 $\Psi[A_z]$：选取 $A_z$（而非 $A_{\bar z}$）作为自变量是这一整段约定的一部分，抄写时切勿把 $z$ 与 $\bar z$ 对调。]

\subsubsection{Weyl 费米子的反常} % 原文: 4.4.2 Anomalies Of Weyl Fermions

WZ 项概括了量子作用量的反常变分。让我们在二维情形下把这一点简要勾勒一下。

设 $\psi_+$ 是 $1+1$ 维中的一个复手征费米子场。我们先用初等量子场论来解释这个反常。

让我们把该费米子与一个 $U(1)$ 规范场耦合，并定义有效作用量：

\begin{equation}
  e^{-\Gamma(A_-)}:=\int d\psi_+d\bar\psi_+\,e^{-\int\bar\psi_+(\partial_-+A_-)\psi_+}
  \tag{4.100}
  \label{eq:4.100}
\end{equation}

形式上我们有

\begin{equation}
  \Gamma(A_-)=-\log\det(\partial_-+A_-)
  \tag{4.101}
  \label{eq:4.101}
\end{equation}

并且形式上有效作用量在下述变换下是规范不变的（$\epsilon$ 为实）：

\begin{equation}
  A_-\to A_-+id\epsilon(x^-)
  \tag{4.102}
  \label{eq:4.102}
\end{equation}

因为我们可以作补偿变换

\begin{equation}
  \begin{aligned}
    \psi_+ &\to e^{-i\epsilon}\psi_+\\
    \bar\psi_+ &\to e^{i\epsilon}\bar\psi_+
  \end{aligned}
  \tag{4.103}
  \label{eq:4.103}
\end{equation}

以得到作用量的不变性。然而，单圈量子效应破坏了这一形式上的规范不变性。表达式 $\Gamma(A)$ 的非不变性就称为反常（anomaly）。

让我们像（图）13 那样用 Feynman 图展开来推导反常的表达式：也就是说，我们写出

\begin{equation}
  \Gamma(A_-)=-\operatorname{Tr}\operatorname{Log}\Big(1+\frac{1}{\partial_-}A_-\Big)
  \tag{4.104}
  \label{eq:4.104}
\end{equation}

% [Tutor Note: 原 PDF 第 119 页的 Figure 13 为空白占位框，图形内容在原稿中缺失，仅存图注；此处按原图注翻译并保留图号，以免后文图号错位。]
\begin{figure}[htbp]
  \centering
  \MissingFigureNote
  \caption{$\Gamma(A_-)$ 按 $A_-$ 的幂级数展开时所出现的 Feynman 图。（原文 Fig.~13，第 119 页）}
  \label{fig:13}
\end{figure}

% [Tutor Note: 原 PDF 第 119 页的 Figure 14 为空白占位框，图形内容在原稿中缺失，仅存图注；此处按原图注翻译并保留图号，以免后文图号错位。]
\begin{figure}[htbp]
  \centering
  \MissingFigureNote
  \caption{动量 $q$ 流入这个两点图。（原文 Fig.~14，第 119 页）}
  \label{fig:14}
\end{figure}

并按 $A$ 展开。二次项就是图 14 的那个图。费米子的传播子是

\begin{equation}
  \langle\bar\psi_+(p)\psi_+(-p)\rangle=\frac{1}{p_-+i\epsilon\operatorname{sign}(p_+)}
  \tag{4.105}
  \label{eq:4.105}
\end{equation}

注意 $i\epsilon$ 项上那个微妙的符号。这可以从标量传播子的 $m\to0$ 极限来理解

\begin{equation}
  \frac{1}{p^2+m^2+i\epsilon}\Big|_{m\to0}=\frac{1}{p_+\big(p_-+i\epsilon\operatorname{sign}(p_+)\big)}
  \tag{4.106}
  \label{eq:4.106}
\end{equation}

于是这个 Feynman 图是

\begin{equation}
  \int dp_+dp_-\,\frac{1}{(p_-+i\epsilon\operatorname{sign}(p_+))(p_--q_-+i\epsilon\operatorname{sign}(p_+-q_+))}
  \tag{4.107}
  \label{eq:4.107}
\end{equation}

注意，如果两个 $i\epsilon$ 项的符号同为正，那么我们可以把 $p_-$ 的积分在另半平面闭合而得到零。因此，为使结果非零，它们的符号必须相反。设 $q_+>0$，则必须有 $0<p_+<q_+$，此时积分的计算是直接的，给出 $q_+/q_-$。这样我们得到

\begin{equation}
  \Gamma(A_-)\sim\int d^2q\,A_-(q)A_-(-q)\frac{q_+}{q_-}+\mathcal{O}(A^3)
  \tag{4.108}
  \label{eq:4.108}
\end{equation}

% [Tutor Note: (4.105) 与 (4.106) 的 $\epsilon\operatorname{sign}(p_+)$ 是全篇最容易抄错的地方：符号函数里的变量是 $p_+$（光锥「时间」分量），不是 $p_-$；而 (4.107) 的分母第二个因子中的符号函数是 $\operatorname{sign}(p_+-q_+)$，正是这个位移使得两个 $i\epsilon$ 在 $0<p_+<q_+$ 区间内符号相反、围道不能同时在同侧闭合。核对 (4.108) 的 $q_+/q_-$ 时，请以 $q_+>0$ 的假设为前提。]

由此得到的高阶图结果是规范不变的，于是在规范变换

\begin{equation}
  A_-(q)\to A_-(q)+iq_-\epsilon(q)
  \tag{4.109}
  \label{eq:4.109}
\end{equation}

之下我们有

\begin{equation}
  \delta\Gamma(A_-)\sim\int d^2q\,(iq_+A_-)\epsilon(-q)=\int d^2q\,F(q)\epsilon(-q)=\int d^2x\,F(x)\epsilon(x)
  \tag{4.110}
  \label{eq:4.110}
\end{equation}

这就是反常。

现在让我们考虑非阿贝尔情形：有一组费米子场 $\psi^i_+$，它在 $SU(n)$ 的 $\mathbf{n}$ 中变换。那么，与外部 $SU(n)$ 规范场耦合时，有效作用量成为

\begin{equation}
  e^{-\Gamma(A_-)}:=\int d\psi_+d\bar\psi_+\,e^{-\int\bar\psi_{+i}\left(\delta^i{}_j\partial_-+(A_-)^i{}_j\right)\psi_+^j}
  \tag{4.111}
  \label{eq:4.111}
\end{equation}

形式上路径积分在下述变换下不变

\begin{equation}
  \begin{aligned}
    \psi &\to g^{-1}\psi\\
    \bar\psi &\to\bar\psi g\\
    A &\to g^{-1}Ag+g^{-1}dg
  \end{aligned}
  \tag{4.112}
  \label{eq:4.112}
\end{equation}

然而，与上面类似的计算给出：

\begin{equation}
  \Gamma(A^g_-)-\Gamma(A_-)=\frac{i}{2\pi}\int\operatorname{Tr}(\partial_+gg^{-1})A_-+F(g)
  \tag{4.113}
  \label{eq:4.113}
\end{equation}

其中 $F(g)$ 不依赖于 $A_-$。

{\small\itshape ［原文标记：需要解释 $F(g)$ 就是 WZ 项：］}
% [原文待办标记，原文 p.120；其上下各有一行星号分隔线，为作者分隔草稿段落所用]

现在，在二维中我们可以对某个群元素 $h$ 把 $A_-$ 写成 $A_-=h^{-1}\partial_-h$，于是我们断定

\begin{equation}
  \Gamma(A_-)\sim iS^+_{WZW}(h)
  \tag{4.114}
  \label{eq:4.114}
\end{equation}

因此，WZW Lagrange 量把一个经典作用量概括了量子反常的效应。

{\small\itshape ［原文标记：当 $\partial_+A_-\neq0$ 时路径依赖性如何？规范化 WZW 作用量；解析延拓的问题］}
% [原文待办标记，原文 p.120；原文为三行大写体的草稿提示，前后各有一行星号分隔线]

\subsection{圈群的最高权表示} % 原文: 4.5 Highest Weight Representations Of Loop Groups

\begin{keytakeaways}
  \item 出发点：目标空间群 $G$ 紧致时，能量有下界的表示就是最高权表示；Verma 模 $M(\lambda)$ 由满足 $E^\alpha v=0$（$\alpha\in\Phi^+$）、$Hv=\lambda(H)v$ 的循环向量生成，其极大真子模的商给出不可约模 $V(\lambda)$。
  \item 有限维类比：$sl(2)$ 的例子中零向量 $v_{2j+1}=(E^-)^{2j+1}|j\rangle$ 生成极大子模 $M(-2j-2)$，商就是维数为 $2j+1$ 的有限维不可约模；$2j\notin\mathbb{Z}_{\geq0}$ 时 Verma 模本身不可约且无穷维。
  \item 仿射情形的新结构：加入 $L_0$ 后自然 Cartan 子代数为 $(H^i_0,K,L_0)$，$L_0$ 的本征值由 Sugawara 公式定为 $\frac{(\lambda,\lambda+2\rho)}{2(k+h^\vee)}$，而 $H^i_n$ 生成的 Heisenberg 代数保证表示必为无穷维。
  \item 可积性约束：要求每个 $sl(2)_{n,\alpha}$ 子代数的作用都是有限维的，于是 $\lambda$ 必须是支配整权且 $k\ge\lambda\cdot\theta$；对 $\widehat{Lsu(2)}_k$ 即 $0\le j\le k/2$，并由 (4.140) 给出相应零向量。
  \item 几何视角：$G/T$ 是紧 Kähler 流形，其相空间量子化给出的 Hilbert 空间正是 $G$ 的不可约表示（Borel-Weil-Bott 定理）。
\end{keytakeaways}

{\small\itshape ［原文标记：需要对 Borel-Weil-Bott 另设一节，特别是关于 $LG/T$ 的情形。］}
% [原文待办标记，原文 p.120；原文该行上下各有一行星号分隔线]

\subsubsection{\texorpdfstring{$\widehat{L\mathfrak{g}}$}{Lg 带帽} 的可积最高权表示} % 原文: 4.5.1 Integrable highest weight representations of \widehat{Lg}

{\small\itshape ［原文标记：本节没有清楚区分什么时候是在不可约表示里工作、什么时候是在 Verma 模里工作。］}
% [原文待办标记（页边），原文 p.120]

当目标空间群 $G$ 紧致时，我们关心的表示是那些能量可以有下界的表示。由于能量动量张量的 Sugawara 形式，这些表示就是\emph{最高权表示}。

为定义它们，我们引入 \emph{Verma 模}的概念，它是 Lie 代数 $\mathfrak{g}$ 的某一类表示。

选取一个 Cartan 子代数 $\mathfrak{t}\subset\mathfrak{g}$。对任意向量 $\lambda\in\mathfrak{t}^*$（不必在权格中），我们如下构造一个表示 $M(\lambda)$，称为 Verma 模。

选取一个正根与负根的分解。那么，Verma 模 $M(\lambda)$ 是由一个向量 $v$ 生成的循环模，该向量满足

\begin{equation}
  \begin{aligned}
    E^\alpha v &= 0\qquad \forall\alpha\in\Phi^+\\
    Hv &= \lambda(H)v
  \end{aligned}
  \tag{4.115}
  \label{eq:4.115}
\end{equation}

人们通过取 $\prod_i X_iv$ 的张成来构建该模，其中 $X_i$ 是负阶梯算符。于是我们有如下的关键

\begin{theorem}
  % 原文为不分编号的 run-in 定理陈述；LaTeX 按 [section] 自动编号
  $M(\lambda)$ 有一个极大真子模 $S(\lambda)$。商
  \begin{equation}
    V(\lambda):=M(\lambda)/S(\lambda)
    \tag{4.116}
    \label{eq:4.116}
  \end{equation}
  是不可约的。
\end{theorem}

$S(\lambda)$ 中的向量称为「零向量（null vectors）」：$S(\lambda)$ 是一些最高权 Verma 模之和，而那些最高权向量称为本原零向量（primitive null vectors）。

\textbf{例：} $\mathfrak{g}=sl(2)$。我们使用基

\begin{equation}
  \begin{aligned}
    [E^+,E^-] &= H\\
    [H,E^\pm] &= \pm2E^\pm
  \end{aligned}
  \tag{4.117}
  \label{eq:4.117}
\end{equation}

转换到物理学的记号是 $H=2J^3$、$E^\pm=J^\pm=J^1\pm iJ^2$。（采用约定 $[J^a,J^b]=i\epsilon^{abc}J^c$，其中 $a,b,c=1,2,3$ 且 $\epsilon^{123}=+1$。）

最高权向量通常记为 $|j\rangle$，满足

\begin{equation}
  \begin{aligned}
    E^+|j\rangle &= 0\\
    H|j\rangle &= 2j|j\rangle
  \end{aligned}
  \tag{4.118}
  \label{eq:4.118}
\end{equation}

我们也把它记为 $v_0$。

Verma 模由 $v_m:=(E^-)^m|j\rangle$（$m\ge0$）张成。它是无穷维的。一个简单的计算表明

\begin{equation}
  H(E^-)^m|j\rangle=(2j-2m)(E^-)^m|j\rangle
  \tag{4.119}
  \label{eq:4.119}
\end{equation}

让我们寻找一个最高权子模。我们需要计算：

\begin{equation}
  \begin{aligned}
    E^+(E^-)^m|j\rangle &= \sum_{i=0}^{m-1}(E^-)^{m-1-i}[E^+,E^-](E^-)^i|j\rangle\\
    &= m\Big(2j-(m-1)\Big)(E^-)^{m-1}|j\rangle
  \end{aligned}
  \tag{4.120}
  \label{eq:4.120}
\end{equation}

因此，存在非平凡子模当且仅当 $2j\in\mathbb{Z}_{\geq0}$。若 $2j\notin\mathbb{Z}_{\geq0}$，则 $M(2j)=V(j)$ 是一个无穷维不可约表示。若 $2j\in\mathbb{Z}_{\geq0}$，则

\begin{equation}
  v_{2j+1}:=(E^-)^{2j+1}|j\rangle
  \tag{4.121}
  \label{eq:4.121}
\end{equation}

生成极大子模，它同构于 $M(-2j-2)$。商

\begin{equation}
  V(j):=M(2j)/M(-2j-2)
  \tag{4.122}
  \label{eq:4.122}
\end{equation}

同构于维数为 $2j+1$ 的有限维不可约模。

\textbf{评注：}

\begin{enumerate}
  \item 通过把 $E^\pm=J^1\pm iJ^2$ 看作物理学的算符，我们可以把这些表示看成 $su(2)$ 的表示。于是我们可以在 Verma 模上引入一个 Hermite 形式，使得 $J^-=(J^+)^\dagger$。由上面所说我们看到
  \begin{equation}
    \langle j|(J^+)^{2j+1}(J^-)^{2j+1}|j\rangle=0
    \tag{4.123}
    \label{eq:4.123}
  \end{equation}
  另一方面，这正是范数平方
  \begin{equation}
    \|\,(J^-)^{2j+1}|j\rangle\,\|^2
    \tag{4.124}
    \label{eq:4.124}
  \end{equation}
  因此在一个幺正表示中它必须为零。由于这个原因，向量 $v_{2j+1}$ 被称为\emph{零向量}。
  \item 只有对那些 $2j+1$ 为正整数的表示，Lie 代数的表示才能指数化（exponentiate）为 Lie 群的表示。事实上还有更多。$SU(2)/U(1)=SL(2,\mathbb{C})/B=\mathbb{C}P^1$ 是一个复流形，而该表示自然地等同于 $SU(2)/U(1)$ 上某个全纯线丛的全纯截面。这是几何量子化（geometric quantization）的一个漂亮例子：对紧致单 Lie 群 $G$ 及其极大环面 $T$，旗流形 $G/T$ 是紧 Kähler 流形，从而是辛流形，而相空间的量子化给出的 Hilbert 空间正是 $G$ 的不可约表示。这就是 Borel-Weil-Bott 定理。一种量子力学的进路使用相空间路径积分 $\int\langle\lambda,g^{-1}\dot g\rangle-H(g)$。该路径积分可以用局部化（localization）精确算出，从而给出 Kirillov 特征标定理与 Weyl 特征标公式。关于这一点的介绍见：
  \begin{enumerate}
    \item A. Alekseev, L. D. Faddeev, Samson L. Shatashvili, ``Quantization of symplectic orbits of compact Lie groups by means of the functional integral,'' J. Geom. Phys. 5:391-406, 1988.
    \item Richard J. Szabo, \emph{Equivariant localization of path integrals}. e-Print: hep-th/9608068
  \end{enumerate}
\end{enumerate}

{\small\itshape ［原文标记：上面部分材料将移到附录 B.3。］}
% [原文待办标记，原文 p.122；其上方有一行星号分隔线]

现在，对于以有限维半单 Lie 代数 $\mathfrak{g}$ 为基础的仿射 Lie 代数 $\widehat{L\mathfrak{g}}$，人们可以展开关于 Cartan 子代数、根与权的整套理论。结果表明，把 $L_0$ 一并加入代数 $\widehat{L\mathfrak{g}}$ 中非常有用。于是，选取 $\mathfrak{g}$ 的一个 Cartan-Weyl 分解后，所得代数的一个自然 Cartan 子代数是

\begin{equation}
  H^i_0,\ K,\ L_0
  \tag{4.125}
  \label{eq:4.125}
\end{equation}

$L_0$ 的加入使人们可以在根空间上定义一个非退化的内积。人们再一次可以构造 Verma 模：生成一个 Verma 模的最高权向量满足

\begin{equation}
  H^i_0|v\rangle=\lambda^i|v\rangle
  \tag{4.126}
  \label{eq:4.126}
\end{equation}

而由于 Sugawara 构造，

\begin{equation}
  L_0|v\rangle=\frac{(\lambda,\lambda+2\rho)}{2(k+h^\vee)}|v\rangle
  \tag{4.127}
  \label{eq:4.127}
\end{equation}

于是最高权条件是

\begin{equation}
  E^\alpha_0|v\rangle=0\qquad \alpha\in\Phi^+
  \tag{4.128}
  \label{eq:4.128}
\end{equation}

\begin{equation}
  J^a_n|v\rangle=0\qquad n>0,\ \forall a
  \tag{4.129}
  \label{eq:4.129}
\end{equation}

{\small\itshape ［原文标记：请统一使用 $v$、$|v\rangle$ 或 $|\lambda\rangle$ 中的一种记号；下文这里全都乱了。］}
% [原文待办标记（页边），原文 p.123]

人们再一次可以构造一个 Verma 模，然后取商以得到不可约模。

人们容易看出，不可约表示将是无穷维的，因为仅由 $H^i_n$ 生成的子代数就已经是一个 Heisenberg 代数

\begin{equation}
  [H^i_n,H^j_m]=kn\delta^{ij}\delta_{n+m,0}
  \tag{4.130}
  \label{eq:4.130}
\end{equation}

而以 $H^i_{-n}$（$n>0$）作用于真空不会产生任何零向量。

% [Tutor Note: 「容易看出」省掉的是：对 (4.130) 的 Heisenberg 代数，$n>0$ 的 $H^i_{-n}$ 是纯粹的产生算符、$H^i_n$ 是湮灭算符，两者对易子只是中心 c-数 $kn\delta^{ij}$，故 Fock 空间上没有收缩关系、不可能产生范数为零的向量。而这恰恰是反面对照：正因为 $E^\alpha_n$ 之间不像 Heisenberg 代数那样平凡，才会出现零向量。]

在这一点上，人们把注意力限制到\emph{可积的}最高权表示上。为构造这些表示，我们注意到对 $\mathfrak{g}$ 的每个根 $\alpha$ 与每个整数 $n$，都可以构造一个由 $J^+\sim E^\alpha_n$、$J^-\sim E^{-\alpha}_{-n}$ 生成的 $sl(2)$ 子代数。

为简单起见，设 $\mathfrak{g}$ 为 A-D-E 型，于是我们可以把根归一化为 $\alpha^2=2$。那么

\begin{equation}
  [E^\alpha_n,E^{-\alpha}_{-n}]=\alpha^iH^i_0+nK
  \tag{4.131}
  \label{eq:4.131}
\end{equation}

于是我们与 $sl(2)$ 的标准基有如下同构：

\begin{equation}
  \begin{aligned}
    H &\to \alpha^iH^i_0+nK\\
    E^+ &\to E^\alpha_n\\
    E^- &\to E^{-\alpha}_{-n}
  \end{aligned}
  \tag{4.132}
  \label{eq:4.132}
\end{equation}

我们称它为 $sl(2)_{n,\alpha}$。

% [Tutor Note: (4.131) 中 $\alpha^iH^i_0$ 里指标 $i$ 是求和指标，$\alpha^i$ 是根 $\alpha$ 在 Cartan 基 $H^i$ 上的分量；千万不要把它误读成 $\alpha$ 的 $i$ 次幂。另外，$K$ 前方的系数与 $H$ 中 $J^3_0$ 的系数是绑定的：写成 $\alpha^iH^i_0+nK$ 时对应 (4.141) 的 $k-2J^3_0$，把 $n\to-n$（等价地整体翻转 $K$ 的中心扩张方向）则对应 (4.142) 的 $k+2J^3_0$；此时 $k$ 与 $h^\vee$ 同时变号，而 (4.133) 的约束形式不变。]

\textbf{定义：} \emph{可积}最高权表示 $V(\lambda)$ 是指这样的表示：$V(\lambda)$ 可以分解为 $sl(2)$ 子代数 $sl(2)_{n,\alpha}$ 的\emph{有限维}表示之和。

换句话说，一旦我们把 Cartan 子代数 $(H^i_0,K,L_0)$ 对角化以得到权向量，每个具有确定权向量的态在 $sl(2)_{n,\alpha}$ 作用下都必须生成一个有限维表示。

{\small\itshape ［原文标记：这里需要解释一些术语并提供更多细节］}
% [原文待办标记（页边），原文 p.124]

对 $sl(2)_{0,\alpha}$ 考虑这一点，我们发现 $\lambda$ 必须是 $\mathfrak{g}$ 的支配整权（dominant integral weight）。然后我们得知，对 $n\ge0$ 与所有 $\alpha$ 有 $nk+\alpha\cdot\lambda\in\mathbb{Z}_{\geq0}$。最强的约束来自 $n=1$ 与 $\alpha=-\theta$，其中 $\theta$ 是最高根。所以 $k$ 必须是非负整数，而可积表示的权受到如下约束

\begin{equation}
  k\ge\lambda\cdot\theta
  \tag{4.133}
  \label{eq:4.133}
\end{equation}

在物理学中，这常常由\emph{幺正性}的要求导出。基本的计算是：

\begin{equation}
  \begin{aligned}
    0 &\le \|\,E^\alpha_{-1}|\lambda\rangle\,\|^2\\
    &= \langle\lambda|E^{-\alpha}_1E^\alpha_{-1}|\lambda\rangle\\
    &= \langle\lambda|[E^{-\alpha}_1,E^\alpha_{-1}]|\lambda\rangle\\
    &= \big(k-(\alpha,\lambda)\big)\,\|\,|\lambda\rangle\,\|^2\\
    &= \big[(k-(\theta,\lambda))+(\theta-\alpha,\lambda)\big]\,\|\,|\lambda\rangle\,\|^2
  \end{aligned}
  \tag{4.134}
  \label{eq:4.134}
\end{equation}

{\small\itshape ［原文标记：此处 $v=|v\rangle=|\lambda\rangle$。请统一记号！］}
% [原文待办标记（页边），原文 p.124]

这里 $\theta$ 是最高根。因此，只要 $\lambda$ 是支配权且 $k-(\theta,\lambda)\ge0$，幺正性约束就得到满足。

让我们考虑 $\widehat{Lsu(2)}$ 的情形。让我们选取一个由下式定义的最高权向量 $|j\rangle$：

\begin{equation}
  \begin{aligned}
    J^3_0|j\rangle &= j|j\rangle\\
    K|j\rangle &= k|j\rangle\\
    J^+_0|j\rangle &= 0\\
    J^a_n|j\rangle &= 0\qquad n>0,\ \forall a
  \end{aligned}
  \tag{4.135}
  \label{eq:4.135}
\end{equation}

由 Sugawara 公式我们算出

\begin{equation}
  L_0|j\rangle=\frac{j(j+1)}{k+2}|j\rangle
  \tag{4.136}
  \label{eq:4.136}
\end{equation}

我们首先用 $J^-_0$ 作用，以填满 $sl(2)_{0,+}$ 的一个自旋 $j$ 多重态：

\begin{equation}
  |j\rangle,\ J^-_0|j\rangle,\ (J^-_0)^2|j\rangle,\ \dots,\ (J^-_0)^{2j}|j\rangle
  \tag{4.137}
  \label{eq:4.137}
\end{equation}

% [Tutor Note: 原 PDF 第 125 页的 Figure 15 为空白占位框，图形内容在原稿中缺失，仅存图注；此处按原图注翻译并保留图号，以免后文图号错位。]
\begin{figure}[htbp]
  \centering
  \MissingFigureNote
  \caption{$\widehat{Lsu(2)}_k$ 的最高权表示中重数非零的点的示意图。（原文 Fig.~15，第 125 页）}
  \label{fig:15}
\end{figure}

% [Tutor Note: 原 PDF 第 125 页的 Figure 16 为空白占位框，图形内容在原稿中缺失，仅存图注；此处按原图注翻译并保留图号，以免后文图号错位。]
\begin{figure}[htbp]
  \centering
  \MissingFigureNote
  \caption{$k=1$、$j=0$ 与 $k=3$、$j=0$ 时 $\widehat{Lsu(2)}_k$ 最高权表示的简并度。取自 I. Affleck 等，J. Phys. A Math. Gen. \textbf{22}(1989)511。（原文 Fig.~16，第 125 页）}
  \label{fig:16}
\end{figure}

注意这些态都有相同的 $L_0$ 本征值。

接下来我们用 $J^+_{-1}$ 作用，得到最高权向量 $J^+_{-1}|j\rangle$。现在我们可以多次用 $J^-_0$ 作用。一般地，我们将得到 Clebsch-Gordon 级数 $(j)\otimes(1)=(j+1)\oplus(j-1)$ 中的那些态。

另一方面，我们也可以让由下式生成的 $sl(2)_{1,-}$ 作用

\begin{equation}
  \begin{aligned}
    H &\to k-2J^3_0\\
    E^+ &\to J^-_1\\
    E^- &\to J^+_{-1}
  \end{aligned}
  \tag{4.138}
  \label{eq:4.138}
\end{equation}

由于对这个代数有 $H|j\rangle=(k-2j)|j\rangle$，我们得到自旋为 $k/2-j$ 的表示，从而得知对一个可积表示：

\begin{equation}
  0\le j\le\frac{k}{2}
  \tag{4.139}
  \label{eq:4.139}
\end{equation}

并且此外

\begin{equation}
  (J^+_{-1})^{k+1-2j}|j\rangle\sim0
  \tag{4.140}
  \label{eq:4.140}
\end{equation}

是一个零向量。

\medskip\noindent\textbf{练习}\quad 证明

\begin{equation}
  \begin{aligned}
    H &\to k-2J^3_0\\
    E^+ &\to J^-_1\\
    E^- &\to J^+_{-1}
  \end{aligned}
  \tag{4.141}
  \label{eq:4.141}
\end{equation}

\begin{equation}
  \begin{aligned}
    H &\to k+2J^3_0\\
    E^+ &\to J^+_1\\
    E^- &\to J^-_{-1}
  \end{aligned}
  \tag{4.142}
  \label{eq:4.142}
\end{equation}

二者都是与 $sl(2)$ 的同构。描述与第二个代数相关的一些零向量。

\medskip\noindent\textbf{练习}\quad 对 $k=1$ 的仿射 $su(2)$，算出其最高权表示的最初几个能级。

{\small\itshape ［原文标记：上一节还需要大幅改进：列出零向量的完整结构——先手工构造一些，然后施加仿射 Weyl 群。这样就能得到此情形下的 BGG 消解。然后引用一般的 BGG 消解。好的参考：Pressley-Segal, \emph{Loop Groups}。］}
% [原文待办标记，原文 p.126；原文该段前后各有一行星号分隔线]

零向量：注意

\begin{equation}
  (E^-_0)^{2j+1}|j\rangle
  \tag{4.143}
  \label{eq:4.143}
\end{equation}

是有限维表示论中的标准结果，其自旋为 $2J^3_0=-j-1$，所以生成那种类型的 Verma 模。

现在

\begin{equation}
  (E^+_{-1})^{k+1-2j}|j\rangle
  \tag{4.144}
  \label{eq:4.144}
\end{equation}

被 $E^+_0$ 湮灭，但具有自旋

\begin{equation}
  2J^3_0=(2k+2)-2j
  \tag{4.145}
  \label{eq:4.145}
\end{equation}

（由对 {\small\itshape ［原文标记：\ensuremath{\ast}\ensuremath{\ast}\ensuremath{\ast}\ensuremath{\ast}\ensuremath{\ast}\ensuremath{\ast}\ensuremath{\ast}］} 作翻转 $j\to(k+2)-(j+1)$ 得到）
% [原文待填空白，原文 p.126；原文此处为一串星号，作者尚未填入具体对象]

于是，现在我们考虑

\begin{equation}
  (E^-_0)^{2k-2j+1}(E^+_{-1})^{k+1-2j}|j\rangle
  \tag{4.146}
  \label{eq:4.146}
\end{equation}

其自旋为 $2J^3_0=$ {\small\itshape ［原文标记：待计算］}，如此等等。
% [原文待填空白，原文 p.126；原文此处为一长串星号，自旋值尚未填出；其后另有两行星号分隔线，为该节草稿的结束标记]
